% Part: incompleteness
% Chapter: incompleteness-provability
% Section: first-incompleteness-thm

\documentclass[../../../include/open-logic-section]{subfiles}

\begin{document}

\olfileid{inc}{inp}{1in}

\olsection{قضیهٔ نخست ناتمامیت}

اکنون می‌توانیم اثبات اصلی گودل برای قضیهٔ نخست ناتمامیت را شرح دهیم.
بگذارید~$\Th{T}$ نظریه‌ای دلخواه و به‌طور محاسبه‌پذیر اصل‌موضوعی‌شده
در زبانی گسترش‌دهندهٔ زبان حساب باشد، به‌گونه‌ای که~$\Th{T}$ اصول
موضوعهٔ~$\Th{Q}$ را در بر گیرد. این امر، به‌ویژه، بدین معناست که
$\Th{T}$ توابع و روابط محاسبه‌پذیر را بازنمایی می‌کند.

استدلال کردیم که با داشتن رمزگذاری معقولی از فرمول‌ها و اثبات‌ها به
صورت اعداد، رابطهٔ~$\Prf[\Th{T}](x,y)$ محاسبه‌پذیر است؛ رابطهٔ
$\Prf[\Th{T}](x,y)$ برقرار است اگر و تنها اگر~$x$ عدد گودل
!!a{derivation}ی از
!!{formula}ی با عدد گودل~$y$ در~$\Th{T}$ باشد. در واقع، گودل برای
نظریهٔ خاص مورد نظر خود توانست با استفاده از فهرست ۴۵ تابع و رابطه در
مقاله‌اش نشان دهد این رابطه بازگشتی اولیه است. رابطهٔ چهل‌وپنجم،
$x B y$، همان~$\Prf[\Th{T}](x,y)$ است، برای انتخاب خاص او از~$\Th{T}$.
به یاد داشته باشید هرجا گودل در مقاله‌اش واژهٔ «بازگشتی» را به
کار می‌برد، امروزه عبارت «بازگشتی اولیه» را به کار می‌بریم.

چون~$\Prf[\Th{T}](x,y)$ محاسبه‌پذیر است، در~$\Th{T}$ بازنمایی‌پذیر
است. از~$\OPrf[\Th{T}](x,y)$ برای اشاره به فرمول بازنمایندهٔ آن استفاده
می‌کنیم. بگذارید~$\OProv[\Th{T}](y)$ فرمول
$\lexists[x][\OPrf[\Th{T}](x,y)]$ باشد. این فرمول رابطهٔ چهل‌وششم،
$\fn{Bew}(y)$، را در فهرست گودل توصیف می‌کند. چنان‌که گودل اشاره
می‌کند، این تنها رابطه‌ای است که «نمی‌توان ادعا کرد بازگشتی است». منظور
محتمل او چنین است: از تعریف روشن نیست که این رابطه محاسبه‌پذیر باشد؛
و تحولات بعدی در واقع نشان می‌دهند که محاسبه‌پذیر نیست.

بگذارید~$\Th{T}$ نظریه‌ای !!{axiomatizable} و شامل~$\Th{Q}$ باشد.
آنگاه~$\Prf[\Th{T}](x, y)$ تصمیم‌پذیر، و ازاین‌رو در~$\Th{Q}$ با
!!a{formula}ی چون~$\OPrf[\Th{T}](x, y)$ بازنمایی‌پذیر
است. بگذارید~$\OProv[\Th{T}](y)$ همان فرمولی باشد که در بالا توصیف
کردیم. بنا بر لم نقطهٔ ثابت، فرمولی چون~$!G_\Th{T}$ وجود دارد که
$\Th{Q}$ (و بنابراین~$\Th{T}$) رابطهٔ زیر را !!{derive}s می‌کند:
\begin{equation}
\ollabel{eqn:qpf}
!G_\Th{T} \liff \lnot \OProv[\Th{T}](\gn{!G_\Th{T}}).
\end{equation}
توجه کنید که~$!G_\Th{T}$ در اصل می‌گوید: «$!G_\Th{T}$ در~$\Th{T}$
!!{derivable} نیست.»

\begin{lem}\ollabel{lem:cons-G-unprov}
اگر~$\Th{T}$ نظریه‌ای سازگار و !!{axiomatizable} و گسترشی از~$\Th{Q}$
باشد، آنگاه~$\Th{T} \Proves/ !G_\Th{T}$.
\end{lem}

\begin{proof}
فرض کنید~$\Th{T}$، $!G_\Th{T}$ را !!{derive}s می‌کند. در این صورت \emph{واقعاً}
!!a{derivation} وجود دارد و بنابراین برای عددی چون~$m$ رابطهٔ
$\Prf[\Th{T}](m, \Gn{!G_\Th{T}})$ برقرار است. اما آنگاه~$\Th{Q}$
جملهٔ~$\OPrf[\Th{T}](\num m, \gn{!G_\Th{T}})$ را !!{derive}s می‌کند. پس
$\Th{Q}$، $\lexists[x][\OPrf[\Th{T}](x,\gn{!G_\Th{T}})]$ را
!!{derive}s می‌کند، که بنا بر تعریف همان
$\OProv[\Th{T}](\gn{!G_\Th{T}})$ است. بنا بر~\olref{eqn:qpf}،
$\Th{Q}$، $\lnot !G_\Th{T}$ را !!{derive}s می‌کند و چون~$\Th{T}$ گسترشی از
$\Th{Q}$ است، $\Th{T}$ نیز چنین می‌کند. نشان دادیم اگر~$\Th{T}$،
$!G_\Th{T}$ را !!{derive}s می‌کند، آنگاه~$\lnot !G_\Th{T}$ را نیز
!!{derive}s می‌کند و در نتیجه ناسازگار خواهد بود.
\end{proof}

\begin{defn}
\ollabel{thm:oconsis-q}
نظریهٔ~$\Th{T}$ را \emph{$\omega$-سازگار} می‌نامیم هرگاه شرط زیر
برقرار باشد: اگر~$\lexists[x][!A(x)]$ جمله‌ای دلخواه باشد و~$\Th{T}$
عبارات~$\lnot !A(\num 0)$، $\lnot !A(\num 1)$،
$\lnot !A(\num 2)$، \dots را !!{derive}s می‌کند، آنگاه~$\Th{T}$
$\lexists[x][!A(x)]$ را ثابت نمی‌کند.
\end{defn}

توجه کنید هر نظریهٔ $\omega$-سازگار، سازگار نیز هست. این امر به‌سادگی
از این واقعیت نتیجه می‌شود که اگر~$\Th{T}$ ناسازگار باشد، آنگاه
$\Th{T} \Proves !A$ برای هر~$!A$. به‌ویژه، اگر~$\Th{T}$ ناسازگار
باشد، هم عبارت~$\lnot !A(\num n)$ را برای هر~$n$ !!{derive}s می‌کند و هم عبارت
$\lexists[x][!A(x)]$ را !!{derive}s می‌کند. پس اگر~$\Th{T}$ ناسازگار باشد،
$\omega$-ناسازگار است. بنا بر عکس نقیض، اگر~$\Th{T}$
$\omega$-سازگار باشد، باید سازگار باشد.

\begin{lem}\ollabel{lem:omega-cons-G-unref}
اگر~$\Th{T}$ نظریه‌ای $\omega$-سازگار و !!{axiomatizable} و گسترشی
از~$\Th{Q}$ باشد، آنگاه~$\Th{T} \Proves/ \lnot !G_\Th{T}$.
\end{lem}

\begin{proof}
نشان می‌دهیم اگر~$\Th{T}$، $\lnot !G_\Th{T}$ را !!{derive}s می‌کند، آنگاه
$\omega$-ناسازگار است. فرض کنید~$\Th{T}$، $\lnot !G_\Th{T}$ را
!!{derive}s می‌کند. اگر~$\Th{T}$ ناسازگار باشد، $\omega$-ناسازگار است و کار
تمام است. در غیر این صورت~$\Th{T}$ سازگار است؛ پس~$!G_\Th{T}$ را
!!{derive} نمی‌کند، بنا بر~\olref{lem:cons-G-unprov}. چون هیچ
!!{derivation}ی از~$!G_\Th{T}$ در~$\Th{T}$ وجود ندارد، $\Th{Q}$
عبارات زیر را !!{derive}s می‌کند:
\[
\lnot \OPrf[\Th{T}](\num 0, \gn{!G_\Th{T}}), \lnot \OPrf[\Th{T}](\num 1,
\gn{!G_\Th{T}}), \lnot \OPrf[\Th{T}](\num 2, \gn{!G_\Th{T}}), \dots
\]
و~$\Th{T}$ نیز چنین می‌کند. از سوی دیگر، بنا بر~\olref{eqn:qpf}،
$\lnot !G_\Th{T}$ هم‌ارز است با
$\lexists[x][\OPrf[\Th{T}](x,\gn{!G_\Th{T}})]$. پس~$\Th{T}$
$\omega$-ناسازگار است.
\end{proof}

\begin{prob}
  هر نظریهٔ $\omega$-سازگار، سازگار است. نشان دهید عکس این گزاره برقرار
  نیست؛ یعنی نظریه‌های سازگار اما $\omega$-ناسازگار وجود دارند. این کار
  را با نشان‌دادن اینکه~$\Th{Q} \cup
  \{\lnot !G_\Th{Q}\}$ سازگار اما $\omega$-ناسازگار است انجام دهید.
\end{prob}

\begin{thm}
\ollabel{thm:first-incompleteness} فرض کنید~$\Th{T}$ نظریه‌ای دلخواه،
$\omega$-سازگار و !!{axiomatizable} و گسترشی از~$\Th{Q}$ باشد. آنگاه
$\Th{T}$ کامل نیست.
\end{thm}

\begin{proof}
  اگر~$\Th{T}$، $\omega$-سازگار باشد، سازگار است؛ پس
  $\Th{T} \Proves/ !G_\Th{T}$، بنا بر~\olref{lem:cons-G-unprov}، برقرار
  است. بنا بر
  \olref{lem:omega-cons-G-unref} نیز
  $\Th{T} \Proves/ \lnot !G_\Th{T}$. این بدان معناست که~$\Th{T}$
  ناتمام است، زیرا نه~$!G_\Th{T}$ و نه~$\lnot !G_\Th{T}$ را !!{derive}s می‌کند.
\end{proof}

\end{document}
